% O template chama esta seção de "Desenvolvimento"; as Diretrizes pedem
% "3. Metodologia" com "3.1 Modelagem" e "3.2 Método de Solução". Seguimos as Diretrizes.
\section{Metodologia}
\label{metodologia}

% ORÇAMENTO: até 2 páginas para 3.1 + 3.2 juntas.
% Escrever PRIMEIRO (é o que trava o código).

\subsection{Modelagem}
\label{ssec:modelagem}

% [Características do dataset] (exigido explicitamente pelas Diretrizes)
% Tabelas usadas, tamanho da amostra (ex.: pedidos "train"), categoria "missing",
% ausência de quantidade, de loja e de layout, order_dow sem mapeamento documentado.

% [Mapeamento categoria -> corredor físico] regra + variações de layout testadas.

% [Modelo 1: grafo da loja] G = (V, E, w) com a notação dos slides da disciplina.
% V: pontas dos corredores e caixa; E: corredores e trechos de travessa; w: metros.
% Corredor INDIVISÍVEL (o dataset não dá a posição do produto): corredor com item = aresta obrigatória.
% Rota mínima = problema do Carteiro Rural (ligação com Euler).
% Custo da simplificação: superestima a distância quando o item está perto da entrada (vai para limitações).
% Figura sugerida (largura total, 2 colunas):
% \begin{figure*}[!ht]
%   \centering
%   \includegraphics[width=0.9\textwidth]{figs/grafo-loja.pdf}
%   \caption{Modelo em grafo da loja ... (o quê, como, onde, quando).}
%   \label{fig-grafoLoja}
% \end{figure*}

% [Modelo 2: grafo de coocorrência] vértices = categorias; arestas = Jaccard; por que não contagem bruta.
% Regra do layout por coocorrência; construído com uma parte dos pedidos e avaliado com outra (evita circularidade).

\subsection{Método de Solução}
\label{ssec:metodo}

% [Visão geral] diagrama do pipeline: dados -> mapeamento -> grafo do pedido -> estratégias -> métricas.

% [Estratégias] tabela em LaTeX (o template proíbe tabela como imagem):
% carrinho, aleatória, vizinho mais próximo (próximo corredor obrigatório), S-shape, ótimo.
% Dizer por que retorno e maior lacuna ficaram de fora: dependem da posição do item no corredor.
% Ótimo: ordem + sentido de travessia de cada corredor obrigatório, ligados por caminhos mínimos
% (programação dinâmica estilo Held-Karp); custo dobra a cada corredor distinto.

% [Algoritmo do núcleo] pseudocódigo com o ambiente algorithm (algorithm2e já carregado pela classe):
% \begin{algorithm}[ht!]
%   \KwData{...}
%   ...
%   \caption{...}
%   \label{alg-...}
% \end{algorithm}

% [Desenho experimental] comparação pareada, estratificação pelo nº de corredores distintos, sementes fixas na ordem aleatória,
% layouts: padrão, variações e coocorrência;
% divisão de pedidos (construir layout x avaliar), análise de sensibilidade.

% [Métricas e parâmetros] distância, tempo (velocidade, tempo por item), pedidos por turno.

% [Implementação] linguagem, bibliotecas, link do repositório GitHub (obrigatório no relatório).
